% Figure from MATH 590 notes, section 7. Compile with the notes preamble.
\begin{tikzpicture}[scale=1.6]
  % X = {a,b,c} with open sets {b}, {a,b}, {b,c}, X, and the empty set
  \draw[black!55, thin, rounded corners=6pt] (-0.5,-0.78) rectangle (1.5,0.9);
  \node[font=\scriptsize, black!60] at (1.36,0.76) {$X$};
  \draw[figB, thick, rounded corners=9pt] (-0.28,-0.3) rectangle (0.7,0.45);
  \draw[figB, thick, rounded corners=9pt] (0.3,-0.45) rectangle (1.28,0.3);
  \draw[figR, thick] (0.52,0.0) circle (0.22);
  \foreach \p/\q in {0/a,0.5/b,1/c} {
    \fill (\p,0) circle (1.1pt);
    \node[font=\scriptsize, below=0.5pt] at (\p,0) {$\q$};
  }
  \node[font=\scriptsize, figB] at (-0.02,0.6) {$\{a,b\}$};
  \node[font=\scriptsize, figB] at (1.02,-0.6) {$\{b,c\}$};
  \draw[figR, thin] (0.6,0.2) -- (0.72,0.64);
  \node[font=\scriptsize, figR, above] at (0.74,0.6) {$\{b\}$};
\end{tikzpicture}
